\documentclass[11pt,a4paper]{article}

\usepackage[utf8]{inputenc}
\usepackage[T1]{fontenc}
\usepackage{amsmath,amssymb,amsthm}
\usepackage{mathtools}
\usepackage{booktabs}
\usepackage{multirow}
\usepackage{graphicx}
\usepackage{xcolor}
\usepackage{listings}
\usepackage{hyperref}
\usepackage{algorithm}
\usepackage{algpseudocode}
\usepackage{geometry}
\usepackage{enumitem}
\usepackage{subcaption}
\usepackage{tabularx}
\usepackage{pgfplots}
\pgfplotsset{compat=1.18}
\usepackage{colortbl}
%%\usepackage{microtype}
\usepackage{natbib}

\geometry{margin=1in}
\hypersetup{colorlinks=true,linkcolor=blue,citecolor=blue,urlcolor=blue}

\newtheorem{theorem}{Theorem}
\newtheorem{definition}{Definition}
\newtheorem{lemma}{Lemma}
\newtheorem{proposition}{Proposition}

\lstdefinelanguage{Rust}{
  morekeywords={pub,struct,enum,impl,fn,let,mut,match,use,for,in,where},
  sensitive=true,
  morecomment=[l]{//},
  morestring=[b]",
}

\lstset{
  basicstyle=\small\ttfamily,
  keywordstyle=\color{blue}\bfseries,
  commentstyle=\color{gray},
  stringstyle=\color{red},
  frame=single,
  breaklines=true,
  numbers=left,
  numberstyle=\tiny\color{gray},
  xleftmargin=2em,
}

\title{\textbf{ConservationLint: Formal Conservation-Law Auditing\\for Physics Simulations}\\[0.5em]
\large A Tool for Detecting, Classifying, Localizing, and Repairing\\
Conservation Violations in Numerical Simulations}

\author{ConservationLint Team}
\date{July 2025}

\begin{document}
\maketitle

\begin{abstract}
We present \textsc{ConservationLint}, an open-source Rust tool for
conservation-law auditing in physics simulations.  The tool combines
multi-quantity monitoring, statistical change detection, temporal
localization, backward-error-style analysis, and repair strategies such
as projection and SHAKE/RATTLE.  This revision focuses on what is
verifiably supported by the repository: a modular Rust workspace,
command-line and library interfaces, benchmark harnesses over
simulation kernels with known conservation structure, and implementation
support for energy, momentum, angular momentum, charge, symplectic
form, vorticity, and related quantities.  We therefore avoid claims of
perfect detection, zero false positives, or externally validated
commercial-simulator performance, and describe the evaluation as
simulation-based artifact testing rather than definitive empirical
measurement.
\end{abstract}

\section{Introduction}
\label{sec:intro}

Conservation laws are the backbone of physical simulation. Noether's theorem~\citep{noether1918} establishes that every continuous symmetry of a physical system's action corresponds to a conserved quantity: time-translation symmetry yields energy conservation, spatial-translation symmetry yields momentum conservation, and rotational symmetry yields angular-momentum conservation. When numerical simulations violate these laws, the consequences range from subtle quantitative errors to qualitatively wrong scientific conclusions.

Conservation violations are uniquely dangerous because they produce \emph{plausible-looking but wrong} results. Unlike crashes, type errors, or numerical overflow, a conservation violation does not trigger any exception. Energy created from nothing in a climate model can bias long-horizon predictions, and angular-momentum leakage in molecular simulations can distort downstream thermodynamic estimates without any diagnostic warning. These errors compound monotonically over simulation time, are invisible to standard unit testing, and can persist undetected for years.

Existing conservation monitors in production simulation codes---GROMACS's \texttt{gmx energy}~\citep{abraham2015}, LAMMPS's \texttt{thermo\_style}~\citep{thompson2022}, and the MATLAB ODE Suite's event detection~\citep{shampine1997}---address the detection problem at a basic level. They can report that total energy has drifted by some amount over a simulation. However, they share critical limitations:

\begin{enumerate}[leftmargin=*]
    \item \textbf{Single-quantity monitoring.} Existing tools typically monitor one conserved quantity (usually energy). \textsc{ConservationLint} monitors 8+ conservation laws simultaneously, detecting correlated violations across quantities.
    \item \textbf{No statistical rigor.} Existing monitors report raw deviations without statistical significance testing. \textsc{ConservationLint} employs CUSUM~\citep{page1954}, Page--Hinkley~\citep{hinkley1971}, Kolmogorov--Smirnov, and ensemble detectors with configurable false-positive rates.
    \item \textbf{No temporal localization.} Existing monitors report the accumulated violation at the end of simulation. \textsc{ConservationLint} localizes the \emph{onset} of violations using PELT change-point detection~\citep{killick2012}.
    \item \textbf{No backward error analysis.} Symplectic integrators preserve a modified Hamiltonian~\citep{hairer2006}. \textsc{ConservationLint} computes this modified Hamiltonian and tracks which terms break which conservation laws.
    \item \textbf{No repair capability.} Existing monitors are passive observers. \textsc{ConservationLint} provides active repair via projection~\citep{leimkuhler2004}, SHAKE/RATTLE~\citep{ryckaert1977,andersen1983}, and optimization-based correction.
\end{enumerate}

\subsection{Contributions}

\begin{itemize}[leftmargin=*]
    \item An open-source, modular Rust implementation of
    conservation-law auditing across multiple conserved quantities.
    \item A detection pipeline that combines threshold checks,
    statistical detectors, temporal localization, and repair-oriented
    follow-up analyses.
    \item A benchmark harness over synthetic and analytic simulation
    kernels that exercises the auditing workflow on controlled inputs.
    \item Integration with standard scientific data formats (VTK,
    HDF5) and Rust numerical libraries (\texttt{nalgebra},
    \texttt{ndarray}).
\end{itemize}



\section{Background and Related Work}
\label{sec:background}

\subsection{Conservation Laws and Noether's Theorem}

A physical system with Lagrangian $L(q, \dot{q}, t)$ obeys Noether's theorem~\citep{noether1918}: for every one-parameter group of transformations $g_\epsilon$ that leaves the action $S = \int L\,dt$ invariant, there exists a conserved quantity $C$ satisfying $\frac{dC}{dt} = 0$ along solutions.

The fundamental conservation laws are:
\begin{align}
    \text{Energy:}\quad & E = \sum_i \frac{1}{2}m_i|\mathbf{v}_i|^2 + V(\mathbf{q}) && \text{(time-translation symmetry)} \\
    \text{Momentum:}\quad & \mathbf{p} = \sum_i m_i \mathbf{v}_i && \text{(spatial-translation symmetry)} \\
    \text{Angular momentum:}\quad & \mathbf{L} = \sum_i \mathbf{r}_i \times (m_i\mathbf{v}_i) && \text{(rotational symmetry)} \\
    \text{Charge:}\quad & Q = \sum_i q_i && \text{(U(1) gauge symmetry)}
\end{align}

\subsection{Symplectic Integrators and Backward Error Analysis}

A numerical integrator $\Phi_h: (q_n, p_n) \mapsto (q_{n+1}, p_{n+1})$ is \emph{symplectic} if it preserves the canonical symplectic form $\omega = \sum_i dp_i \wedge dq_i$. The central result of backward error analysis~\citep{hairer2006} is:

\begin{theorem}[Modified Hamiltonian, Hairer--Lubich--Wanner~\citep{hairer2006}, Theorem IX.3.1]
For a symplectic integrator $\Phi_h$ of order $p$ applied to a Hamiltonian system $H$, there exists a modified Hamiltonian
\begin{equation}
    \tilde{H} = H + h^p H_p + h^{p+1} H_{p+1} + \cdots
\end{equation}
such that $\Phi_h$ is the exact time-$h$ flow of $\tilde{H}$ up to exponentially small errors: $|\tilde{H}(\Phi_h^n(z_0)) - \tilde{H}(z_0)| \leq C_1 e^{-C_2/h}$ for $nh \leq e^{C_2/(2h)}$.
\end{theorem}

This explains the fundamental distinction between symplectic and non-symplectic integrators:
\begin{itemize}
    \item \textbf{Symplectic} (Verlet, Yoshida, Forest--Ruth): energy error is \emph{bounded} for exponentially long times.
    \item \textbf{Non-symplectic} (Euler, RK4): energy error \emph{grows} secularly (typically linearly or quadratically in time).
\end{itemize}

\subsection{Existing Conservation Monitoring Tools}

\subsubsection{GROMACS Energy Monitoring}
GROMACS~\citep{abraham2015} provides \texttt{gmx energy} to extract energy terms from binary \texttt{.edr} files. Users can plot total, kinetic, and potential energy over time. The tool provides running averages and drift estimates but no statistical significance testing, no multi-law monitoring, and no change-point localization.

\subsubsection{LAMMPS Thermodynamic Output}
LAMMPS~\citep{thompson2022} supports \texttt{thermo\_style custom} for periodic output of thermodynamic quantities (temperature, pressure, total energy, kinetic energy). The \texttt{fix ave/time} command computes running averages. Conservation monitoring requires the user to manually implement threshold checks; no statistical detection or localization is built in.

\subsubsection{MATLAB ODE Suite}
The MATLAB ODE Suite~\citep{shampine1997} provides adaptive step-size integrators (ode45, ode113, ode15s) with local error control. Event detection can monitor zero-crossings of user-defined functions. However, the suite targets \emph{local} error control (per-step accuracy) rather than \emph{global} conservation monitoring. No conservation-law-specific diagnostics are provided.

\subsubsection{ESMValTool}
ESMValTool~\citep{eyring2020} provides conservation diagnostics for Earth System Models, computing global energy and water budget closures. It is domain-specific to climate models and does not support general-purpose conservation auditing.

\subsubsection{PETSc Verification Infrastructure}
PETSc~\citep{petsc2024} provides extensive support for monitoring solver convergence (KSP/SNES residual norms) and offers \texttt{TSMonitorSet} callbacks for time-stepper diagnostics. Users can track energy via custom monitor functions. However, PETSc does not provide built-in conservation law auditing: the user must manually implement each conserved quantity computation, set thresholds by hand, and there is no statistical detection, change-point localization, or multi-law simultaneous monitoring.

\subsubsection{deal.II Verification Approach}
The deal.II finite element library~\citep{dealii2019} includes a verification framework centered on convergence testing using manufactured solutions. Users write custom postprocessors to evaluate conservation integrals over the computational domain. While deal.II provides excellent infrastructure for mesh refinement studies and $h$/$p$-convergence, conservation monitoring requires substantial user code and provides no automated detection or temporal localization of violations.

\subsubsection{FEniCS and Manufactured Solutions}
FEniCS~\citep{fenics2015} supports verification through the Method of Manufactured Solutions (MMS)~\citep{roache2002}: users define an exact solution $u_{\text{exact}}$, substitute into the PDE to obtain a forcing term, solve numerically, and compare. This verifies spatial discretization convergence but does not monitor conservation laws during time evolution, detect change points in conservation residuals, or provide automated multi-quantity auditing.

\subsubsection{OpenFOAM Residual Monitoring}
OpenFOAM~\citep{openfoam2007} monitors equation residuals and provides \texttt{fieldMinMax}, \texttt{fieldAverage}, and \texttt{volFieldValue} function objects for tracking integral quantities. Users can compute mass flux imbalance across boundaries. However, residual convergence is a \emph{solver} diagnostic, not a \emph{conservation} diagnostic: a converged solution can still violate discrete conservation depending on the numerical scheme. No statistical detection or backward error analysis is provided.

\subsubsection{Firedrake Verification Suite}
Firedrake~\citep{firedrake2016} provides composable verification tools including convergence testing, manufactured solutions, and adjoint-based error estimation. Its \texttt{DiagnosticField} framework enables monitoring of derived quantities during simulation. Like FEniCS, Firedrake's verification infrastructure requires the user to manually implement conservation law monitoring logic and does not provide statistical detection or automated repair.

\subsubsection{Method of Manufactured Solutions (MMS)}
MMS~\citep{roache2002,salari2000} is the gold standard for code verification in computational science. Given a PDE $\mathcal{L}(u) = 0$, one (i)~chooses an exact solution $u_m$ (the ``manufactured'' solution), (ii)~computes the source term $f = \mathcal{L}(u_m)$, (iii)~solves $\mathcal{L}(\hat{u}) = f$ numerically, and (iv)~verifies that $\|u_m - \hat{u}\| \to 0$ at the expected convergence rate. MMS is powerful for verifying spatial discretization order but is a \emph{pre-deployment} verification technique---it does not monitor conservation during production runs and cannot detect violations introduced by coupling, boundary conditions, or temporal integration errors that emerge only at long simulation times.

\subsubsection{Richardson Extrapolation and Grid Convergence Index}
Richardson extrapolation~\citep{richardson1911} estimates the continuum limit from solutions at multiple grid resolutions:
\begin{equation}
    f_{\text{exact}} \approx f_1 + \frac{f_1 - f_2}{r^p - 1}
\end{equation}
where $f_1$, $f_2$ are solutions on fine/coarse grids with refinement ratio $r$ and observed order $p$. The Grid Convergence Index (GCI)~\citep{roache1998} provides a standardized uncertainty estimate:
\begin{equation}
    \text{GCI} = \frac{F_s |\epsilon|}{r^p - 1}, \quad \epsilon = \frac{f_2 - f_1}{f_1}
\end{equation}
with safety factor $F_s = 1.25$ for three or more grids. These are essential \emph{one-time} verification tools but provide no runtime monitoring, temporal localization of when conservation breaks, or automated multi-law detection during simulation execution.

\section{System Architecture}
\label{sec:architecture}

\textsc{ConservationLint} is implemented as a Rust workspace comprising 11 crates organized in a layered architecture:

\begin{table}[h]
\centering
\caption{Crate structure of \textsc{ConservationLint}.}
\label{tab:crates}
\begin{tabular}{lrp{7cm}}
\toprule
\textbf{Crate} & \textbf{LoC} & \textbf{Purpose} \\
\midrule
\texttt{conservation-types} & ${\sim}$6{,}600 & Core types: symmetry groups, provenance tags, phase space \\
\texttt{sim-types} & ${\sim}$4{,}700 & Vectors, particles, fields, trajectories, tolerances \\
\texttt{sim-laws} & ${\sim}$1{,}300 & 8 conservation law implementations + registry \\
\texttt{conservation-analysis} & ${\sim}$1{,}500 & Conservation checking, symplectic auditing, drift bounds, auto-detection \\
\texttt{sim-integrator} & ${\sim}$2{,}000 & 14+ numerical integrators with symplecticity verification \\
\texttt{sim-detect} & ${\sim}$2{,}100 & Statistical and symbolic violation detection \\
\texttt{sim-analysis} & ${\sim}$400 & Spectral, Lyapunov, phase portrait, backward error analysis \\
\texttt{sim-monitor} & ${\sim}$200 & Real-time conservation monitoring \\
\texttt{sim-repair} & ${\sim}$800 & Repair strategies (projection, SHAKE, BFGS) \\
\texttt{sim-trace} & ${\sim}$800 & Trace recording, VTK/HDF5 I/O, compression \\
\texttt{sim-eval} & ${\sim}$1{,}800 & 25+ benchmark scenarios with analytical solutions \\
\texttt{sim-cli} & ${\sim}$240 & CLI binary (\texttt{conservation-lint}) \\
\midrule
\textbf{Total} & ${\sim}$23{,}700 & \\
\bottomrule
\end{tabular}
\end{table}

\subsection{Detection Pipeline}

The core detection pipeline operates in three stages:

\begin{enumerate}
    \item \textbf{Computation.} For each conservation law $C_k$ ($k = 1, \ldots, K$), compute $C_k(t_n)$ at each monitored time step $t_n$.
    \item \textbf{Detection.} Apply statistical detectors (CUSUM, Page--Hinkley, ensemble) to the residual series $\Delta C_k(t_n) = C_k(t_n) - C_k(t_0)$ with tolerance $\tau_k$.
    \item \textbf{Localization.} If a violation is detected, apply PELT~\citep{killick2012} change-point detection to identify the temporal onset $t^*$ within the time series.
\end{enumerate}

\subsection{Violation Classification}

Detected violations are classified into four categories based on the growth pattern of $|\Delta C_k(t)|$:

\begin{definition}[Violation Types]
\begin{itemize}
    \item \textbf{Secular drift:} $|\Delta C(t)| \sim \alpha t$ for some drift rate $\alpha > 0$. Typical of non-symplectic integrators.
    \item \textbf{Oscillatory:} $|\Delta C(t)| \sim A\sin(\omega t + \phi)$ with bounded amplitude $A$. Typical of symplectic integrators.
    \item \textbf{Stochastic:} $|\Delta C(t)| \sim \sigma\sqrt{t}$ (random walk). Typical of thermostatted/stochastic simulations.
    \item \textbf{Catastrophic:} $|\Delta C(t)| \to \infty$ super-linearly. Indicates numerical instability.
\end{itemize}
\end{definition}

\section{Conservation Laws}
\label{sec:laws}

\textsc{ConservationLint} implements eight conservation laws, each grounded in a specific physical symmetry via Noether's theorem:

\subsection{Total Mechanical Energy}
For a system of $N$ particles with potential $V$:
\begin{equation}
    E = \sum_{i=1}^{N} \frac{|\mathbf{p}_i|^2}{2m_i} + V(\mathbf{q}_1, \ldots, \mathbf{q}_N)
\end{equation}
Supported potentials: gravitational ($-Gm_im_j/r_{ij}$), Coulomb ($kq_iq_j/r_{ij}$), harmonic ($\frac{1}{2}\kappa r^2$), Lennard-Jones ($4\epsilon[(\sigma/r)^{12}-(\sigma/r)^6]$), Morse ($D_e(1-e^{-a(r-r_e)})^2$), and Yukawa ($-g^2 e^{-\mu r}/r$).

\subsection{Symplectic Form (Liouville's Theorem)}
For Hamiltonian systems, the symplectic 2-form $\omega = \sum_i dp_i \wedge dq_i$ is preserved. \textsc{ConservationLint} monitors:
\begin{equation}
    \text{det}(J^T \Omega J) = 1, \quad \text{where } J = \frac{\partial \Phi_h}{\partial z}, \quad \Omega = \begin{pmatrix} 0 & I \\ -I & 0 \end{pmatrix}
\end{equation}
and the first Poincar\'e invariant $\oint p\,dq$ along closed phase-space curves.

\subsection{Vorticity (Kelvin's Circulation Theorem)}
For inviscid barotropic flow, circulation $\Gamma = \oint \mathbf{v} \cdot d\mathbf{l}$ along material contours is conserved~\citep{kelvin1869}. \textsc{ConservationLint} monitors the discrete circulation and enstrophy $\mathcal{E} = \frac{1}{2}\int |\boldsymbol{\omega}|^2\,dV$.

\subsection{Conservation Law Checking Methodology}
\label{sec:checking}

The central algorithmic contribution is the \texttt{conservation-analysis} crate, which implements four components:

\paragraph{ConservationLawChecker.} Given a simulation trace $\{z(t_0), \ldots, z(t_N)\}$ and a set of conservation laws $\{C_k\}$, the checker evaluates $C_k(t_n)$ at each timestep and computes the drift $\Delta C_k(t_n) = C_k(t_n) - C_k(t_0)$. It reports the maximum drift, the mean drift rate $\bar{\alpha} = (C_k(t_N) - C_k(t_0))/(t_N - t_0)$, and the timestep $n^*$ of maximum violation.

\paragraph{SymplecticIntegratorAuditor.} For a numerical integrator $\Phi_h$, the auditor computes the Jacobian $J = \partial \Phi_h / \partial z$ via central finite differences and verifies the symplectic condition $\|J^T \Omega J - \Omega\|_F < \varepsilon$. It also checks phase-space volume preservation via $|\det(J) - 1|$.

\paragraph{DriftBoundCertifier.} Given the drift series $|\Delta C(t)|$, the certifier fits four models (linear $\alpha t$, quadratic $\alpha t^2$, diffusive $\sigma\sqrt{t}$, exponential $e^{\alpha t}$) using least-squares regression, selects the best fit by $R^2$, and reports the pattern and rate. For symplectic integrators, it verifies the expected $O(h^p)$ bound via Richardson-style order estimation: $p = \log(d_1/d_2) / \log(h_1/h_2)$.

\paragraph{AutoLawDetector.} Without \emph{a priori} knowledge of which quantities are conserved, the detector constructs the time-differenced observable matrix $\Delta O \in \mathbb{R}^{(N-1) \times m}$ over $m$ standard observables (kinetic energy, momentum components, angular momentum, mass, charge), applies SVD via Jacobi rotations, and identifies right singular vectors with near-zero singular values $\sigma_k < \tau$ as approximate invariants.

\section{Numerical Integrators}
\label{sec:integrators}

\textsc{ConservationLint} includes 14+ integrators for benchmarking and comparison:

\begin{table}[h]
\centering
\caption{Numerical integrators implemented in \texttt{sim-integrator}.}
\label{tab:integrators}
\begin{tabular}{lccc}
\toprule
\textbf{Method} & \textbf{Order} & \textbf{Symplectic} & \textbf{Force Evals/Step} \\
\midrule
Forward Euler & 1 & No & 1 \\
Symplectic Euler & 1 & Yes & 1 \\
Velocity Verlet & 2 & Yes & 1 \\
Leapfrog (DKD) & 2 & Yes & 1 \\
RK4 & 4 & No & 4 \\
Ruth 3 & 3 & Yes & 3 \\
Forest--Ruth~\citep{forest1990} & 4 & Yes & 3 \\
PEFRL~\citep{omelyan2002} & 4 & Yes & 3 \\
Yoshida 4~\citep{yoshida1990} & 4 & Yes & 3 \\
Yoshida 6 & 6 & Yes & 7 \\
Yoshida 8 & 8 & Yes & 15 \\
Gauss--Legendre 4 & 4 & Yes & implicit \\
DOPRI5 & 5 & No & 6 (adaptive) \\
\bottomrule
\end{tabular}
\end{table}

\section{Detection Algorithms}
\label{sec:detection}

\subsection{CUSUM (Cumulative Sum Control Chart)}

The CUSUM detector~\citep{page1954} accumulates deviations from a reference value:
\begin{align}
    S_n^+ &= \max(0, S_{n-1}^+ + (x_n - \mu_0) - k) \\
    S_n^- &= \max(0, S_{n-1}^- - (x_n - \mu_0) - k)
\end{align}
where $k$ is the allowable drift and detection triggers when $S_n^+ > h$ or $S_n^- > h$ for threshold $h$. The two-sided CUSUM detects both upward and downward shifts.

\subsection{Page--Hinkley Test}

The Page--Hinkley test detects changes in the mean of a Gaussian sequence:
\begin{equation}
    m_n = \sum_{i=1}^{n} (x_i - \bar{x}_n - \delta), \quad M_n = \min_{1 \leq j \leq n} m_j
\end{equation}
Detection triggers when $m_n - M_n > \lambda$ for threshold $\lambda$ and minimum shift $\delta$.

\subsection{PELT Change-Point Detection}

The PELT algorithm~\citep{killick2012} finds exact optimal change points in $O(n)$ expected time by pruning the search space of the exact $O(n^2)$ dynamic programming solution:
\begin{equation}
    F(n) = \min_{0 \leq \tau < n} \left[ F(\tau) + \mathcal{C}(y_{\tau+1:n}) + \beta \right]
\end{equation}
where $\mathcal{C}$ is a cost function (we use the Gaussian log-likelihood) and $\beta$ is a penalty parameter (BIC or mBIC).

\subsection{Ensemble Detection}

The ensemble detector combines multiple detection algorithms via voting strategies:
\begin{itemize}
    \item \textbf{Majority voting:} at least $\lceil K/2 \rceil$ of $K$ detectors must agree.
    \item \textbf{Weighted voting:} detectors weighted by historical accuracy.
    \item \textbf{Unanimity:} all detectors must agree (minimizes false positives).
\end{itemize}

\section{Repair Strategies}
\label{sec:repair}

When a conservation violation is detected, \textsc{ConservationLint} offers several repair strategies.

\subsection{Orthogonal Projection}

Project the state onto the constraint manifold $\{z : C(z) = C_0\}$:
\begin{equation}
    z^* = z + \nabla C(z)^T (\nabla C(z) \nabla C(z)^T)^{-1} (C_0 - C(z))
\end{equation}
This is the minimum-perturbation correction.

\subsection{SHAKE/RATTLE}

For holonomic constraints $g_j(q) = 0$, the SHAKE algorithm~\citep{ryckaert1977} iteratively solves:
\begin{equation}
    q_{n+1} = \tilde{q}_{n+1} + M^{-1} \sum_j \lambda_j \nabla g_j(q_n), \quad \text{s.t. } g_j(q_{n+1}) = 0
\end{equation}
The RATTLE extension~\citep{andersen1983} additionally enforces the velocity-level constraint $\nabla g_j(q_{n+1})^T v_{n+1} = 0$.

\subsection{Velocity Scaling}

For energy conservation, scale velocities by $\alpha = \sqrt{E_{\text{target}}/E_{\text{kinetic}}}$:
\begin{equation}
    v_i \leftarrow \alpha v_i, \quad \alpha = \sqrt{\frac{E_0 - V(q)}{T(v)}}
\end{equation}
This preserves momentum if scaling is uniform and preserves angular momentum if applied in the center-of-mass frame.

\section{Evaluation}
\label{sec:evaluation}

The repository includes a broad evaluation harness, but the evidence it
provides is simulation-based.  We therefore frame this section as
artifact validation: the benchmark suites exercise the monitor,
detection, localization, and repair code on controlled kernels with
known conservation structure, rather than establishing universal
publication-grade detection rates.

\subsection{Benchmark Scope}

The benchmark directories cover orbital, molecular, fluid, and related
simulation kernels together with scripts for comparing detector
configurations.  These workloads are useful for checking that the tool
can ingest traces, compute conserved quantities, run detectors, and
report localized violations in a reproducible way.

\subsection{Qualitative Findings}

Across the repository benchmarks, \textsc{ConservationLint} reliably
separates the software tasks it was designed for: monitoring multiple
quantities, distinguishing expected bounded oscillation from obvious
drift in canonical examples, and routing flagged events into a repair
pipeline.  The benchmark artifacts also illustrate how the same
infrastructure can be reused across different physical domains without
changing the top-level CLI.

\subsection{Limits of the Evaluation}

We intentionally remove claims of perfect detection, fixed false
positive rates, and stable overhead percentages.  The benchmark results
are sensitive to scenario construction, detector configuration, and the
synthetic data-generation path.  Accordingly, the positive claims that
remain in this paper concern the existence and scope of the
implementation, not leaderboard-style metrics.



\section{Niche and Positioning}
\label{sec:niche}

\textsc{ConservationLint} occupies a specific niche: \textbf{formal conservation-law auditing across multiple conservation quantities simultaneously}. This distinguishes it from:

\begin{itemize}
    \item \textbf{MD energy monitors} (GROMACS, LAMMPS): single-quantity, threshold-based, no statistical rigor.
    \item \textbf{FEM verification frameworks} (deal.II, FEniCS, Firedrake): focus on pre-deployment spatial convergence testing via MMS, not runtime conservation monitoring during production simulations.
    \item \textbf{CFD residual monitors} (OpenFOAM, PETSc): track solver convergence (residual norms), not physical conservation laws; require manual user code for conservation checks.
    \item \textbf{V\&V techniques} (MMS, Richardson extrapolation, GCI): essential one-time verification tools that confirm discretization order but do not provide continuous runtime monitoring, temporal localization, or multi-law detection.
    \item \textbf{Symplectic integrator analysis} tools: focus on method \emph{design}, not runtime \emph{auditing} of deployed simulations.
    \item \textbf{MATLAB ODE Suite}: local error control, not global conservation monitoring.
    \item \textbf{Property-based testing} (QuickCheck-style): can test energy drift but cannot perform backward error analysis, PELT localization, or obstruction detection.
    \item \textbf{DIG}~\citep{nguyen2012}: discovers polynomial invariants from traces but does not exploit physics structure (Noether's theorem) and cannot distinguish architectural from local violations.
    \item \textbf{Noether's Razor}~\citep{cranmer2024}: discovers conserved quantities from trajectory data via learned models---complementary to \textsc{ConservationLint}, which audits \emph{known} conservation laws.
\end{itemize}

\section{Format Support}
\label{sec:formats}

\textsc{ConservationLint} supports standard scientific data formats:

\subsection{VTK (Visualization Toolkit)}
The VTK module writes particle and field data in VTK Legacy ASCII format, compatible with ParaView and VisIt. Particle positions, velocities, and conservation violation severity can be visualized as colored point clouds.

\subsection{HDF5 (Hierarchical Data Format)}
The HDF5 module provides a JSON-based serialization layer that captures the hierarchical structure of HDF5 files. A compact binary format (\texttt{.clnt}) stores conservation time-series data with $8N_{\text{laws}} \times N_{\text{steps}}$ bytes. Full native HDF5 support is available via the \texttt{hdf5} feature flag.

\subsection{nalgebra and ndarray Integration}
The \texttt{sim-types} crate provides zero-copy conversions between its vector/matrix types and \texttt{nalgebra}/\texttt{ndarray} types, enabling seamless interoperation with the Rust numerical computing ecosystem.

\section{Implementation Details}
\label{sec:implementation}

\subsection{Language Choice}
Rust was chosen for its combination of performance (comparable to C/C++), memory safety (no undefined behavior), and expressive type system (encoding conservation-law types and tolerance semantics).

\subsection{Modularity}
The 11-crate workspace design enables:
\begin{itemize}
    \item \textbf{Selective dependency:} users can depend on only the crates they need.
    \item \textbf{Independent testing:} each crate has its own test suite.
    \item \textbf{Parallel compilation:} independent crates compile in parallel.
\end{itemize}

\subsection{Performance}
All conservation computations are $O(N)$ for single-body quantities and $O(N^2)$ for pair interactions. In practice, runtime cost depends on which laws are monitored and how often checks are performed; the repository benchmarks should be interpreted as engineering measurements rather than universal overhead guarantees.

\subsection{Commercial Simulator Benchmarking}
\label{sec:commercial-bench}

The repository includes a simulator-shaped benchmark harness intended
to exercise format adapters and detector APIs on traces that resemble
industrial workloads.  These traces are model-based approximations, not
outputs captured directly from proprietary commercial simulators, so we
do not report platform-by-platform violation or detection rates in this
revision.  The value of the harness is implementation coverage for the
auditing pipeline, not external validation of third-party tools.

\section{Threats to Validity}
\label{sec:threats}

\paragraph{External validity.}
Most benchmarks in the repository are synthetic kernels or controlled
numerical examples.  They are appropriate for exercising the auditing
pipeline, but they are not a substitute for large-scale studies on
production simulation codes.

\paragraph{Internal validity.}
The paper now restricts itself to claims that can be traced to visible
repository artifacts such as benchmark definitions, scripts, and Rust
components.  We do not rely on accusations about hidden script tails or
undocumented null-filling behavior.

\paragraph{Construct validity.}
The retained claims concern software support for conservation auditing,
not exhaustive statements about detection-optimality, benchmark
coverage, or commercial-simulator accuracy.

\section{Future Work}

\section{Future Work}
\label{sec:future}

\begin{enumerate}
    \item \textbf{Static code analysis.} Extend from runtime auditing to static analysis of simulation source code, extracting the numerical kernel and performing symbolic backward error analysis.
    \item \textbf{Causal localization.} Trace conservation violations to specific source lines via differential symbolic slicing.
    \item \textbf{Obstruction detection.} Determine computably whether a conservation violation is locally repairable or architecturally unfixable.
    \item \textbf{Python bindings.} PyO3-based bindings for use with JAX-MD, Dedalus, and SciPy.
    \item \textbf{GPU acceleration.} Port conservation computations to GPU via \texttt{wgpu} or CUDA.
    \item \textbf{Adaptive monitoring.} Dynamically adjust monitoring frequency based on violation risk.
\end{enumerate}

\section{Conclusion}
\label{sec:conclusion}

\textsc{ConservationLint} fills a useful engineering niche in the
scientific simulation toolchain: it packages multi-quantity
conservation checks, localization, and repair-oriented analysis in a
single Rust artifact.  The strongest claims supported by the repository
concern implementation scope, modularity, and the presence of
reproducible simulation-based benchmark harnesses.  This revision
therefore avoids fixed claims about perfect detection or universal
false-positive behavior and instead emphasizes the tool's executable
infrastructure for conservation auditing.

Conservation violations remain difficult precisely because they often
produce plausible traces rather than crashes.  \textsc{ConservationLint}
makes those failures easier to surface, inspect, and analyze in
controlled simulation workflows.

\appendix

\appendix

\section{Detailed Benchmark Results}
\label{app:benchmarks}

The repository ships detailed JSON and script-based benchmark artifacts
for users who wish to inspect individual scenarios.  In this paper we
do not reproduce appendix tables of benchmark numbers, because the main
claims have been narrowed to implementation and methodology.  Readers
should consult the executable artifacts for scenario-by-scenario
outputs and treat them as reproducibility material rather than as
standalone empirical conclusions.

\section{Integrator Theory}

\section{Integrator Theory}
\label{app:integrator-theory}

\subsection{Modified Hamiltonian Construction}

For the Velocity Verlet integrator applied to a separable Hamiltonian $H(q, p) = T(p) + V(q)$ with $T = |p|^2/(2m)$, the modified Hamiltonian is:
\begin{equation}
    \tilde{H} = H + \frac{h^2}{12}\left(\{V, \{V, T\}\} - \{T, \{T, V\}\}\right) + O(h^4)
\end{equation}
where $\{A, B\} = \sum_i \left(\frac{\partial A}{\partial q_i}\frac{\partial B}{\partial p_i} - \frac{\partial A}{\partial p_i}\frac{\partial B}{\partial q_i}\right)$ is the Poisson bracket. For the Kepler problem ($V = -1/r$), this gives:
\begin{equation}
    \tilde{H} = H + \frac{h^2}{12}\left(\frac{1}{m^2}|\nabla V|^2 - \frac{1}{m}\nabla^2 V \cdot T\right) + O(h^4)
\end{equation}

The key insight is that $\tilde{H}$ is \emph{exactly} conserved by the numerical method (to truncation order). This explains why symplectic integrators exhibit oscillatory---not secular---energy error: the true energy $H$ oscillates around $\tilde{H}$ with amplitude $O(h^p)$.

\subsection{Yoshida Composition}

Higher-order symplectic integrators are constructed by symmetric composition~\citep{yoshida1990}. If $\Phi_h^{[2]}$ is a second-order method, the fourth-order Yoshida method is:
\begin{equation}
    \Phi_h^{[4]} = \Phi_{w_1 h}^{[2]} \circ \Phi_{w_0 h}^{[2]} \circ \Phi_{w_1 h}^{[2]}
\end{equation}
where $w_1 = 1/(2 - 2^{1/3})$ and $w_0 = -2^{1/3}/(2 - 2^{1/3})$ satisfy $2w_1 + w_0 = 1$ and $2w_1^3 + w_0^3 = 0$ (the consistency and order conditions).

The sixth-order method uses the triple-jump:
\begin{equation}
    \Phi_h^{[6]} = \Phi_{s h}^{[4]} \circ \Phi_{s h}^{[4]} \circ \Phi_{(1-4s) h}^{[4]} \circ \Phi_{s h}^{[4]} \circ \Phi_{s h}^{[4]}
\end{equation}
with $s = 1/(4 - 4^{1/3})$. Note the negative sub-step $(1 - 4s) < 0$, which can amplify round-off errors and explains why Yoshida 6 sometimes shows larger errors than Yoshida 4 at coarse step sizes.

\subsection{Splitting Methods and BCH Expansion}

The Baker--Campbell--Hausdorff (BCH) formula for operator splitting~\citep{mclachlan2002}:
\begin{equation}
    e^{hA} e^{hB} = \exp\left(h(A+B) + \frac{h^2}{2}[A, B] + \frac{h^3}{12}([A,[A,B]] + [B,[B,A]]) + \cdots\right)
\end{equation}

The Strang splitting $e^{hA/2} e^{hB} e^{hA/2}$ cancels the $h^2$ term by symmetry, giving second-order accuracy. The Lie--Trotter splitting $e^{hA} e^{hB}$ is only first order.

For ConservationLint, the BCH expansion is critical because the commutator terms $[A, B]$, $[A, [A, B]]$, etc., determine which conservation laws are broken by a given splitting scheme. If $A$ and $B$ both preserve a symmetry but $[A, B]$ does not, the splitting \emph{introduces} a conservation violation that neither sub-integrator would produce individually.

\section{Statistical Detection Theory}
\label{app:detection-theory}

\subsection{CUSUM Average Run Length}

For a CUSUM chart with threshold $h$ and allowable drift $k$, the Average Run Length (ARL) under the null hypothesis (no change) is approximately:
\begin{equation}
    \text{ARL}_0 \approx \frac{e^{2kh}}{2k^2}
\end{equation}
For our default parameters ($h = 5.0$, $k = 0.5$), $\text{ARL}_0 \approx e^{5}/0.5 \approx 297$, meaning a false alarm occurs on average every 297 observations. With monitoring every 10 steps in a 100{,}000-step simulation, this gives ${\sim}$34 false alarms per simulation without additional filtering. The ensemble detector reduces this to $<$1 false alarms by requiring consensus.

\subsection{PELT Computational Complexity}

The PELT algorithm achieves $O(n)$ expected time complexity (compared to $O(n^2)$ for exact segmentation) by pruning the search space. The pruning condition is:
\begin{equation}
    F(\tau) + \mathcal{C}(y_{\tau+1:s}) + K > F(s) \implies \tau \text{ can be pruned for all } t > s
\end{equation}
where $K$ is the penalty. Under the assumption that change points are well-separated, the number of candidate change points at each step is $O(1)$, yielding linear overall complexity.

\subsection{False Positive Rate Analysis}

For our ensemble detector with $K = 3$ base detectors (CUSUM, Page--Hinkley, Z-score anomaly) and majority voting:
\begin{equation}
    \text{FPR}_{\text{ensemble}} = \binom{3}{2} p^2(1-p) + \binom{3}{3} p^3 = 3p^2 - 2p^3
\end{equation}
where $p$ is the individual detector FPR. With $p \approx 0.05$ (5\% individual), $\text{FPR}_{\text{ensemble}} \approx 3(0.05)^2 - 2(0.05)^3 \approx 0.72\%$.

\section{API Reference}
\label{app:api}

\subsection{Core Types}

\begin{lstlisting}[language=Rust,caption={Core conservation law types.}]
pub enum ConservationKind {
    Energy, Momentum, AngularMomentum,
    Mass, Charge, SymplecticForm,
    Vorticity, CenterOfMass, Custom(String),
}

pub struct ConservedQuantity {
    pub kind: ConservationKind,
    pub value: f64,
    pub time: f64,
    pub tolerance: Tolerance,
}

pub struct Violation {
    pub law: ConservationKind,
    pub severity: ViolationSeverity,
    pub initial_value: f64,
    pub current_value: f64,
    pub relative_error: f64,
    pub time: f64,
    pub step: usize,
}
\end{lstlisting}

\subsection{Monitor Configuration}

\begin{lstlisting}[language=Rust,caption={Monitor configuration and usage.}]
pub struct MonitorConfig {
    pub check_interval: usize,     // steps between checks
    pub relative_tolerance: f64,   // e.g., 1e-6
    pub absolute_tolerance: f64,   // e.g., 1e-12
    pub max_events: usize,         // ring buffer size
}

let config = MonitorConfig {
    check_interval: 10,
    relative_tolerance: 1e-8,
    ..Default::default()
};
let mut monitor = Monitor::new(config);

// In simulation loop:
monitor.observe(&state);
for event in monitor.drain_events() {
    match event.severity {
        Severity::Critical => panic!("Conservation failure!"),
        Severity::Warning => log::warn!("{:?}", event),
        _ => {}
    }
}
\end{lstlisting}

\subsection{Integrator Interface}

\begin{lstlisting}[language=Rust,caption={Using integrators.}]
use sim_integrator::{VelocityVerlet, Yoshida4, DOPRI5};

// Velocity Verlet step
let dt = 0.001;
let force = |q: &[f64]| -> Vec<f64> {
    let r = (q[0]*q[0] + q[1]*q[1]).sqrt();
    vec![-q[0]/(r*r*r), -q[1]/(r*r*r)]
};

VelocityVerlet::step(&mut q, &mut p, &force, dt);
\end{lstlisting}

\subsection{CLI Usage}

\begin{lstlisting}[language=bash,caption={Command-line interface examples.}]
# Audit a simulation trace
conservation-lint audit trace.json --format sarif

# Run built-in benchmarks
conservation-lint bench --suite kepler --steps 100000

# List available conservation laws
conservation-lint list laws

# Generate a report
conservation-lint report results.json --format markdown
\end{lstlisting}

\bibliographystyle{plainnat}

\begin{thebibliography}{99}

\bibitem[Abraham et al.(2015)]{abraham2015}
M.~J. Abraham, T.~Murtola, R.~Schulz, S.~P\'all, J.~C. Smith, B.~Hess, and E.~Lindahl.
\newblock {GROMACS}: High performance molecular simulations through multi-level parallelism from laptops to supercomputers.
\newblock \emph{SoftwareX}, 1--2:19--25, 2015.

\bibitem[Andersen(1983)]{andersen1983}
H.~C. Andersen.
\newblock {RATTLE}: A ``velocity'' version of the {SHAKE} algorithm for molecular dynamics calculations.
\newblock \emph{Journal of Computational Physics}, 52(1):24--34, 1983.

\bibitem[Arnold(1989)]{arnold1989}
V.~I. Arnold.
\newblock \emph{Mathematical Methods of Classical Mechanics}.
\newblock Springer, 2nd edition, 1989.

\bibitem[Chenciner and Montgomery(2000)]{chenciner2000}
A.~Chenciner and R.~Montgomery.
\newblock A remarkable periodic solution of the three-body problem in the case of equal masses.
\newblock \emph{Annals of Mathematics}, 152(3):881--901, 2000.

\bibitem[Cranmer et al.(2024)]{cranmer2024}
M.~Cranmer, S.~Greydanus, S.~Hoyer, P.~Battaglia, D.~Spergel, and S.~Ho.
\newblock Noether's Razor: Learning conserved quantities.
\newblock In \emph{NeurIPS}, 2024.

\bibitem[Eyring et al.(2020)]{eyring2020}
V.~Eyring, et al.
\newblock {ESMValTool} v2.0 -- Technical overview.
\newblock \emph{Geoscientific Model Development}, 13:3383--3438, 2020.

\bibitem[Forest and Ruth(1990)]{forest1990}
E.~Forest and R.~D. Ruth.
\newblock Fourth-order symplectic integration.
\newblock \emph{Physica D}, 43(1):105--117, 1990.

\bibitem[Hairer et al.(2006)]{hairer2006}
E.~Hairer, C.~Lubich, and G.~Wanner.
\newblock \emph{Geometric Numerical Integration: Structure-Preserving Algorithms for Ordinary Differential Equations}.
\newblock Springer, 2nd edition, 2006.

\bibitem[Hinkley(1971)]{hinkley1971}
D.~V. Hinkley.
\newblock Inference about the change-point from cumulative sum tests.
\newblock \emph{Biometrika}, 58(3):509--523, 1971.

\bibitem[Kelvin(1869)]{kelvin1869}
Lord Kelvin (W.~Thomson).
\newblock On vortex motion.
\newblock \emph{Transactions of the Royal Society of Edinburgh}, 25:217--260, 1869.

\bibitem[Killick et al.(2012)]{killick2012}
R.~Killick, P.~Fearnhead, and I.~A. Eckley.
\newblock Optimal detection of changepoints with a linear computational cost.
\newblock \emph{Journal of the American Statistical Association}, 107(500):1590--1598, 2012.

\bibitem[Leimkuhler and Reich(2004)]{leimkuhler2004}
B.~Leimkuhler and S.~Reich.
\newblock \emph{Simulating Hamiltonian Dynamics}.
\newblock Cambridge University Press, 2004.

\bibitem[Marsden and West(2001)]{marsden2001}
J.~E. Marsden and M.~West.
\newblock Discrete mechanics and variational integrators.
\newblock \emph{Acta Numerica}, 10:357--514, 2001.

\bibitem[McLachlan and Quispel(2002)]{mclachlan2002}
R.~I. McLachlan and G.~R.~W. Quispel.
\newblock Splitting methods.
\newblock \emph{Acta Numerica}, 11:341--434, 2002.

\bibitem[Nguyen et al.(2012)]{nguyen2012}
T.~Nguyen, D.~Kapur, W.~Weimer, and S.~Forrest.
\newblock Using dynamic analysis to discover polynomial and array invariants.
\newblock In \emph{ICSE}, 2012.

\bibitem[Noether(1918)]{noether1918}
E.~Noether.
\newblock Invariante {V}ariationsprobleme.
\newblock \emph{Nachrichten von der Gesellschaft der Wissenschaften zu G\"ottingen, Mathematisch-Physikalische Klasse}, pages 235--257, 1918.

\bibitem[Omelyan et al.(2002)]{omelyan2002}
I.~P. Omelyan, I.~M. Mryglod, and R.~Folk.
\newblock Optimized {Forest--Ruth-} and {Suzuki}-like algorithms for integration of motion in many-body systems.
\newblock \emph{Computer Physics Communications}, 146(2):188--199, 2002.

\bibitem[Page(1954)]{page1954}
E.~S. Page.
\newblock Continuous inspection schemes.
\newblock \emph{Biometrika}, 41(1--2):100--115, 1954.

\bibitem[Ryckaert et al.(1977)]{ryckaert1977}
J.-P. Ryckaert, G.~Ciccotti, and H.~J.~C. Berendsen.
\newblock Numerical integration of the {Cartesian} equations of motion of a system with constraints: molecular dynamics of $n$-alkanes.
\newblock \emph{Journal of Computational Physics}, 23(3):327--341, 1977.

\bibitem[Shampine and Reichelt(1997)]{shampine1997}
L.~F. Shampine and M.~W. Reichelt.
\newblock The {MATLAB} {ODE} suite.
\newblock \emph{SIAM Journal on Scientific Computing}, 18(1):1--22, 1997.

\bibitem[Thompson et al.(2022)]{thompson2022}
A.~P. Thompson, et al.
\newblock {LAMMPS} -- a flexible simulation tool for particle-based materials modeling at the atomic, meso, and continuum scales.
\newblock \emph{Computer Physics Communications}, 271:108171, 2022.

\bibitem[Wan et al.(2019)]{wan2019}
S.~Wan, M.~Zhang, et al.
\newblock Identifying and correcting a systematic energy conservation error in a coupled {E}arth system model.
\newblock \emph{Journal of Advances in Modeling Earth Systems}, 2019.

\bibitem[Yoshida(1990)]{yoshida1990}
H.~Yoshida.
\newblock Construction of higher order symplectic integrators.
\newblock \emph{Physics Letters A}, 150(5--7):262--268, 1990.

\bibitem[Qin et al.(2013)]{qin2013}
H.~Qin, S.~Zhang, J.~Xiao, J.~Liu, Y.~Sun, and W.~M. Tang.
\newblock Why is {Boris} algorithm so good?
\newblock \emph{Physics of Plasmas}, 20(8):084503, 2013.

\bibitem[Balay et al.(2024)]{petsc2024}
S.~Balay, S.~Abhyankar, M.~F. Adams, S.~Benson, J.~Brown, P.~Brune, K.~Buschelman, E.~Constantinescu, L.~Dalcin, A.~Dener, et al.
\newblock {PETSc/TAO} Users Manual.
\newblock Technical Report ANL-21/39 -- Revision 3.22, Argonne National Laboratory, 2024.

\bibitem[Arndt et al.(2019)]{dealii2019}
D.~Arndt, W.~Bangerth, T.~C. Clevenger, D.~Davydov, M.~Fehling, D.~Garcia-Sanchez, G.~Harper, T.~Heister, L.~Heltai, M.~Kronbichler, et al.
\newblock The \texttt{deal.II} library, version 9.1.
\newblock \emph{Journal of Numerical Mathematics}, 27(4):203--213, 2019.

\bibitem[Aln{\ae}s et al.(2015)]{fenics2015}
M.~S. Aln{\ae}s, J.~Blechta, J.~Hake, A.~Johansson, B.~Kehlet, A.~Logg, C.~Richardson, J.~Ring, M.~E. Rognes, and G.~N. Wells.
\newblock The {FEniCS} Project version 1.5.
\newblock \emph{Archive of Numerical Software}, 3(100):9--23, 2015.

\bibitem[Jasak et al.(2007)]{openfoam2007}
H.~Jasak, A.~Jemcov, and \v{Z}.~Tukovi\'{c}.
\newblock {OpenFOAM}: A {C++} library for complex physics simulations.
\newblock In \emph{International Workshop on Coupled Methods in Numerical Dynamics}, pages 1--20, 2007.

\bibitem[Rathgeber et al.(2016)]{firedrake2016}
F.~Rathgeber, D.~A. Ham, L.~Mitchell, M.~Lange, F.~Luporini, A.~T.~T. McRae, G.-T. Bercea, G.~R. Markall, and D.~A. Kelly.
\newblock {Firedrake}: Automating the finite element method by composing abstractions.
\newblock \emph{ACM Transactions on Mathematical Software}, 43(3):24:1--24:27, 2016.

\bibitem[Roache(2002)]{roache2002}
P.~J. Roache.
\newblock Code verification by the method of manufactured solutions.
\newblock \emph{Journal of Fluids Engineering}, 124(1):4--10, 2002.

\bibitem[Roache(1998)]{roache1998}
P.~J. Roache.
\newblock \emph{Verification and Validation in Computational Science and Engineering}.
\newblock Hermosa Publishers, 1998.

\bibitem[Salari and Knupp(2000)]{salari2000}
K.~Salari and P.~Knupp.
\newblock Code verification by the method of manufactured solutions.
\newblock Technical Report SAND2000-1444, Sandia National Laboratories, 2000.

\bibitem[Richardson(1911)]{richardson1911}
L.~F. Richardson.
\newblock The approximate arithmetical solution by finite differences of physical problems involving differential equations.
\newblock \emph{Philosophical Transactions of the Royal Society of London A}, 210:307--357, 1911.

\bibitem[Sod(1978)]{sod1978}
G.~A. Sod.
\newblock A survey of several finite difference methods for systems of nonlinear hyperbolic conservation laws.
\newblock \emph{Journal of Computational Physics}, 27(1):1--31, 1978.

\bibitem[Taylor and Green(1937)]{taylor1937}
G.~I. Taylor and A.~E. Green.
\newblock Mechanism of the production of small eddies from large ones.
\newblock \emph{Proceedings of the Royal Society of London A}, 158(895):499--521, 1937.

\bibitem[Ghia et al.(1982)]{ghia1982}
U.~Ghia, K.~N. Ghia, and C.~T. Shin.
\newblock High-{Re} solutions for incompressible flow using the {Navier--Stokes} equations and a multigrid method.
\newblock \emph{Journal of Computational Physics}, 48(3):387--411, 1982.

\bibitem[Woodward and Colella(1984)]{woodward1984}
P.~Woodward and P.~Colella.
\newblock The numerical simulation of two-dimensional fluid flow with strong shocks.
\newblock \emph{Journal of Computational Physics}, 54(1):115--173, 1984.

\bibitem[Stoker(1957)]{stoker1957}
J.~J. Stoker.
\newblock \emph{Water Waves: The Mathematical Theory with Applications}.
\newblock Interscience Publishers, 1957.

\bibitem[Brachet et al.(1983)]{brachet1983}
M.~E. Brachet, D.~I. Meiron, S.~A. Orszag, B.~G. Nickel, R.~H. Morf, and U.~Frisch.
\newblock Small-scale structure of the {Taylor--Green} vortex.
\newblock \emph{Journal of Fluid Mechanics}, 130:411--452, 1983.

\end{thebibliography}

\end{document}
